\subsubsection{Zaglavljene niti i deadlock}\label{gdb-deadlock}

Kod višedretvenih programa simptom često nije pad, već proces koji ne završava. \texttt{gdb} je tada koristan jer možemo zaustaviti proces i pogledati gde čeka svaka nit.

Primer namerno zaključava dva muteksa u različitom redosledu:

\codefile{deadlock.c}{c}

Pokrenimo program iz \texttt{gdb}-a:

\begin{minted}{bash}
$ make
$ gdb ./deadlock
\end{minted}

Program će se najčešće zaglaviti posle kratkog čekanja. U \texttt{gdb}-u ga prekidamo kombinacijom \texttt{Ctrl-C}, zatim pregledamo niti:

\begin{minted}{text}
(gdb) run
^C
(gdb) info threads
(gdb) thread apply all bt
\end{minted}

Komanda \texttt{thread apply all bt} prikazuje stek svake niti. U ovom primeru jedna nit čeka zaključavanje \texttt{right}, a druga čeka zaključavanje \texttt{left}. Pošto svaka već drži drugi muteks, nijedna ne može da nastavi.

Za detaljniju analizu biramo konkretne niti:

\begin{minted}{text}
(gdb) thread 2
(gdb) bt
(gdb) frame 1
(gdb) info locals
(gdb) thread 3
(gdb) bt
\end{minted}

Kod realnih programa stekovi obično vode do viših funkcija aplikacije, ne samo do \texttt{pthread\_mutex\_lock}. Zato se zaključak izvodi iz kombinacije: koji resurs se čeka, koja nit ga drži i kojim redosledom se resursi zaključavaju.

Popravka je dizajnersko pravilo, ne samo promena u debageru: svi delovi programa moraju zaključavati deljene resurse u istom globalnom redosledu ili koristiti viši mehanizam koji taj redosled nameće.
